%======================================================================
\section{The game}\label{sec:game}
%======================================================================

A linear system game sees a group only through its solution group. To turn the Pauli configuration-lamp group into a game we therefore
embed a group that maps onto it into a solution group, with $J$ in the role of the central involution of the game. Slofstra's embedding
theorem does this for groups given by presentations of a special shape, called $E\Gamma_0$~\cite[\S4]{Slo19}: each relation says that two
involutions commute, that some commuting involutions multiply to the identity, or that a generator conjugates one involution to another.
Slofstra and Vidick showed that the embedding keeps a distinguished involution $w$ of the presentation visible in approximate
representations, as the element $J$ of the solution group~\cite[\S6]{SV18}. This section writes down such a presentation $\Pres$ for our
group: its group maps onto $\GPp$ with $w\mapsto J$, and every word of $\RP$ follows from its relations with area at most $60$
(\cref{prop:presentation}). It then proves Theorem~A (\cref{thm:main}), by composing the transfer of \cref{sec:transfer}, the embedding of
\cref{app:gadget}, the sector theorem of \cref{sec:sector} and the hash models of \cref{sec:upper}, once in each direction.

The presentation has only one non-involutive generator, $p$, and this is forced. Houghton's commutator $s=[p,q]$ is an involution, but
in an $E\Gamma_0$ group with two non-involutive generators $p,q$ the relation $[p,q]^2=e$ already implies $[p,q]=e$
(\cref{prop:one-noninv}). We therefore adjoin the ray swap $\iota$ of \cref{lem:GPp} and write $q=\iota p\iota$. Then the commutator
becomes a product of $p$-conjugates of the single involution $\iota$ (equation~\eqref{game:s-lemma}), which the $E\Gamma_0$ format can
express.

What has to be checked is finite. Every relation of $\Pres$ must hold in $\GPp$, so that the hash models of $\GPp$ give approximate
representations of $\Pres$ with $w\mapsto-I$. Every word of $\RP$ must follow from the relations of $\Pres$ with bounded area, so that
approximate representations of $\Pres$ give approximate representations of $\RP$. And Slofstra's construction, applied to $\Pres$, must
produce the linear system of the game. These three facts are verified by machine, each by at least two programs that share no code
(\cref{app:cert}); every derivation can also be replayed by hand.

\subsection{Presentations of type \texorpdfstring{$E\Gamma_0$}{EGamma0}}

\begin{definition}[$E\Gamma_0$ presentations; {\cite[Def.~4.7]{Slo19}, \cite[\S6]{SV18}}]\label{def:EG0}
An \emph{$E\Gamma_0$ presentation} $\Pres=\langle S_0\sqcup S_1\mid R\rangle$ consists of two finite sets of generators, the
\emph{involutive} generators $S_0$ and the linearly ordered \emph{non-involutive} generators $S_1=\{y_1,\dots,y_\ell\}$, a finite family
$\mathcal V$ of subsets of $S_0$ (the \emph{linear sets}), and the following relations, and no others:
\begin{enumerate}[label=(E\arabic*),itemsep=0pt]
\item $x^2=e$ for every $x\in S_0$;
\item for every $V\in\mathcal V$, the \emph{linear relation} $\prod_{x\in V}x=e$, the product taken in a fixed order, and the commutations
$xy=yx$ for $x,y\in V$, which are listed among the relations (E3) as $xyx=y$;
\item conjugacy relations $xyx=z$ for the triples $(x,y,z)$ of a finite set $C_0\subseteq S_0^3$;
\item conjugacy relations $yxy^{-1}=z$ for the triples $(y,x,z)$ of a finite set $C_1\subseteq S_1\times S_0\times S_0$;
\item for some pairs $i>j$, one relation $y_iy_jy_i^{-1}=y_j^{L_{ij}}$ for each such pair, with an integer $L_{ij}\ge1$.
\end{enumerate}
We write $\langle\Pres\rangle$ for the group presented, and $R$ for its relators $x^2$, $\prod_{x\in V}x$, $xyxz^{-1}$, $yxy^{-1}z^{-1}$ and
$y_iy_jy_i^{-1}y_j^{-L_{ij}}$. Areas, approximate representations and state-weighted models of $\Pres$ are taken with respect to $R$.
\end{definition}

Slofstra calls the groups so presented extended homogeneous-linear-plus-conjugacy groups. Because the generators in $S_0$ are
involutions, the relation $xyx=y$ says the same as $xy=yx$, so listing the commutations of a linear set as triples $(x,y,y)\in C_0$, instead of
as the separate relations $xy=yx$ of~\cite[Def.~4.7]{Slo19} and~\cite[\S6]{SV18}, presents the same group. When $\ell\le1$ there are no pairs $i>j$, so (E5) is empty, and if moreover every linear set has at most three
elements, every relator has length at most $4$.

\subsection{Why one non-involutive generator}

In Houghton's group the commutator $s=[p,q]$ is the transposition of $(1,1)$ and $(1,2)$, and the first word of $\RP$ is $s^2$. In an
$E\Gamma_0$ group, the quotient by the normal subgroup generated by the involutive generators sees only the relations (E5), and with two
non-involutive generators the group these define is torsion-free.

\begin{proposition}[Two non-involutive generators are too many]\label{prop:one-noninv}
Let $\Pres$ be an $E\Gamma_0$ presentation with $S_1=\{p,q\}$, and let $G'=\langle\Pres\rangle$. If $[p,q]^2=e$ in $G'$, then $[p,q]=e$
in $G'$. In particular, if $G'$ has a homomorphism to $\GP$ that sends $p$ and $q$ to the generators $p$ and $q$ of $\GP$, then the word
$s^2=[p,q]^2$ of $\RP$ is not trivial in $G'$, so $\RP$ does not follow from the relations of $\Pres$ under the substitution $p\mapsto p$,
$q\mapsto q$.
\end{proposition}

\begin{proof}
Let $N$ be the normal subgroup of $G'$ generated by $S_0$. Every relation of types (E1)--(E4) becomes trivial in $G'/N$, so
$G'/N=\langle p,q\mid R_5\rangle$, where $R_5$ is the set of relations of type (E5). There is at most one of them, since there is one pair
$i>j$. If $R_5$ is empty, $G'/N$ is free of rank $2$. Otherwise $G'/N$ is the Baumslag--Solitar group $\langle a,b\mid bab^{-1}=a^L\rangle$
for some $L\ge1$, which is $\Z[1/L]\rtimes\Z$, and equals $\Z^2$ for $L=1$. These groups are torsion-free. Since $[p,q]^2=e$ in $G'/N$, we
get $[p,q]=e$ in $G'/N$, so $G'/N$ is abelian. A free group of rank two and $\Z[1/L]\rtimes\Z$ with $L\ge2$ are not abelian, so $R_5$
consists of the single relation $qpq^{-1}=p$ or $pqp^{-1}=q$. That relation belongs to the presentation, so $[p,q]=e$ in $G'$.

For the second statement, such a homomorphism sends $[p,q]$ to $s$, which moves the configuration $e_{(1,1)}$ and so is not trivial in
$\GP$. Hence $[p,q]\ne e$ in $G'$, and by the first statement $[p,q]^2\ne e$ in $G'$.
\end{proof}

We keep $p$ as the only non-involutive generator and replace $q$ by $\iota p\iota$, which equals $q$ in $\GPp$ (\cref{lem:GPp}). Write
$\iota_k=p^k\iota p^{-k}$. In the free product $\langle p\rangle*\langle\iota\mid\iota^2\rangle$,
\[
[p,\iota p\iota]=p\iota p\iota p^{-1}\iota p^{-1}\iota=\iota_1\,\iota_2\,\iota_1\,\iota_0 .
\]
So the commutator, which carried the torsion obstruction, is a product of $p$-conjugates of one involution, each reached from $\iota$ by
conjugacy relations $p\iota_kp^{-1}=\iota_{k+1}$ of type (E4); and $\iota_1\iota_2\iota_1$ is again an involution, which a computation in
$Y$ identifies. From the definition of $p$, the involution $\iota_1$ exchanges $(1,1)\leftrightarrow(3,1)$ and $(2,k)\leftrightarrow(3,k+1)$
for $k\ge1$, and fixes $(1,j)$ for $j\ge2$; and $\iota_2$ exchanges $(1,2)\leftrightarrow(3,1)$, $(1,1)\leftrightarrow(3,2)$ and
$(2,k)\leftrightarrow(3,k+2)$ for $k\ge1$, and fixes $(1,j)$ for $j\ge3$. Conjugating $\iota_2$ by $\iota_1$ transports these pairs to
$(1,2)\leftrightarrow(1,1)$, $(3,1)\leftrightarrow(2,1)$ and $(3,k+1)\leftrightarrow(2,k+1)$, and fixes $(1,j)$ for $j\ge3$. Hence
\begin{equation}\label{game:s-lemma}
\iota_1\iota_2\iota_1=\sigma\,\iota_0,\qquad\sigma=\bigl((1,1)\ (1,2)\bigr),
\end{equation}
where $\sigma$ and $\iota_0=\iota$ have disjoint supports; so $[p,\iota p\iota]=\sigma=s$ in $\GPp$. This is the \emph{s-lemma}: one
conjugacy relation $\iota_1\iota_2\iota_1=g_3$ and one linear relation on $\{g_3,\iota,\sigma\}$ express $[p,\iota p\iota]=\sigma$ with
$\sigma\in S_0$. From then on $s$ is an involutive generator, and every other word of $\RP$ says that a conjugate of an involution is an
involution, that two such conjugates commute, or that three commuting ones multiply to $e$, which an $E\Gamma_0$ presentation can express.
The presentation below avoids the obstruction of \cref{prop:one-noninv}: its quotient by the normal subgroup generated by $S_0$ is
$\langle p\rangle\cong\Z$, which has no nontrivial commutators, while by~\eqref{game:s-lemma} the commutator $s$ lies in that subgroup.

\subsection{The presentation \texorpdfstring{$\Pres$}{P}}

\begin{proposition}[The presentation]\label{prop:presentation}
There is an $E\Gamma_0$ presentation $\Pres=\langle S_0\sqcup S_1\mid R\rangle$ with $S_1=\{p\}$ and a distinguished involutive generator
$w$, listed in full in \cref{cert:presentation}, with the following properties.
\begin{enumerate}[label=(\alph*),itemsep=1pt]
\item \emph{Size.} $S_0$ has $26$ elements: $\iota,t,d,x,z$, an involution $\sigma$, and $g_1,\dots,g_{20}$, with $w=g_{14}$. There are $8$
linear sets, each with three elements; $|C_0|=76$, consisting of $32$ triples used by the derivations of (d), the $25$ triples $(g,w,w)$
for $g\in S_0\setminus\{w\}$, and the $19$ commutations inside linear sets not already among these; and $|C_1|=10$. Every relator has
length at most $4$.
\item \emph{A model.} Each $g\in S_0$ comes with a defining word $\vartheta(g)$ in $p,\iota,t,d,x,z$, among them $\vartheta(\iota)=\iota$,
$\vartheta(t)=t$, $\vartheta(d)=d$, $\vartheta(x)=x$, $\vartheta(z)=z$, $\vartheta(\sigma)=[p,\iota p\iota]$, $\vartheta(g_1)=p\iota p^{-1}$,
$\vartheta(g_7)=ptp^{-1}$ and $\vartheta(w)=xzxz$. Put $\vartheta(p)=p$. For every relator $r\in R$ the word $\vartheta(r)$ represents
$1$ in $\GPp$, and its free reduction has length at most $\ell_E=44$. Hence $\vartheta$ induces a homomorphism
$\pi:\langle\Pres\rangle\to\GPp$, which is onto, with $\pi(w)=J$.
\item \emph{Centrality.} $R$ contains $gwg=w$ for every $g\in S_0\setminus\{w\}$, and $pwp^{-1}=w$.
\item \emph{The certificate.} Put $\hat a=a$ for $a\in\{p,t,d,x,z\}$ and $\hat q=\iota p\iota$. For $r\in\RP$ let $\hat r$ be the word in
$p,\iota,t,d,x,z$ obtained by replacing every letter $a^{\pm1}$ of $r$ by $\hat a^{\pm1}$. Then every $\hat r$, $r\in\RP$, and the word
$xzxz\,w^{-1}$ are trivial in $\langle\Pres\rangle$, each with area at most $A_*=60$ with respect to $R$.
\end{enumerate}
The finite statements in (b) and (d) are checked by machine (\cref{app:cert}).
\end{proposition}

The eight linear sets are $\{\iota_1\iota_2\iota_1,\iota,\sigma\}$ (the s-lemma), $\{t,t_2,tt_2\}$, $\{xzx,z,w\}$, $\{zxz,x,w\}$, and the
four sets $\{uvu,v,uvuv\}$ and $\{vuv,u,uvuv\}$ for $(u,v)=(x,t)$ and $(z,t)$, which express the commutation of the Pauli lamp at $0$ with
the one at $e_{(1,1)}$; here $t_2=ptp^{-1}$, and we name generators by their defining words. The ten triples of $C_1$ say that conjugation by
$p$ maps $\iota\mapsto\iota_1\mapsto\iota_2$, $\sigma\mapsto p\sigma p^{-1}\mapsto p^2\sigma p^{-2}$, $t\mapsto t_2\mapsto t_3$ and
$d\mapsto pdp^{-1}$, and fixes $x$, $z$ and $w$.

\begin{proof}
The proposition is a finite list of identities. The programs that check them are described in \cref{app:cert}; here we say what is
checked and why that proves the statement.

\emph{Format and size.} The relations are stored as the lists $C_0$, $C_1$ and $\mathcal V$, so each has one of the shapes (E1)--(E4); the
programs check $S_1=\{p\}$, the counts in (a), and that every pair $u,v$ inside a linear set has $(u,v,v)$ or $(v,u,u)$ in $C_0$. The
length bound is read off the shapes.

\emph{The model, (b).} Each defining word, and each relator translated through $\vartheta$, is evaluated exactly in a faithful model of
$\GPp$ (\cref{cert:model}): every defining word of an involutive generator squares to the identity, every translated relator gives the
identity, and after free reduction the longest translated relator, that of the triple $(g_5,w,w)$, has $44$ letters. So $\vartheta$
induces $\pi$. It is onto because $\GPp$ is generated by $p,\iota,t,d,x,z$ (\cref{lem:GPp}), and $\pi(w)=J$ because $\vartheta(w)=xzxz$.

\emph{Centrality, (c).} The $25$ triples $(g,w,w)$ and the triple $pwp^{-1}=w$ are in the lists.

\emph{The certificate, (d).} This is a statement about the abstract group $\langle\Pres\rangle$ and does not use the model. For each of the
$29$ target words, namely the s-lemma $[p,\iota p\iota]\sigma^{-1}$, the $27$ words $\hat r$ with each occurrence of $[p,\iota p\iota]$
replaced by $\sigma$, and $xzxz\,w^{-1}$, the file records a derivation, which is replayed in three stages.
\begin{enumerate}[label=(\roman*),itemsep=0pt]
\item The word, which has $p$-exponent $0$, is written as a product $\prod_kp^{n_k}g_k^{\pm1}p^{-n_k}$ of conjugates of involutive letters,
$n_k$ being the $p$-exponent of the prefix before $g_k$. This is an identity in the free group, checked by free reduction.
\item Each $p^{n_k}g_kp^{-n_k}$ is replaced by an involutive generator $h_k$ through $\abs{n_k}$ relations of $C_1$, and $g_k^{-1}$ by $g_k$
through (E1). The result is the \emph{resolved sequence} $h_1\cdots h_M$.
\item The resolved sequence is reduced to the empty word by recorded moves, each an instance of a listed relation: cancel $aa$ (E1);
replace $aba$ by $c$ for $(a,b,c)\in C_0$; replace $ab$ by $ba$ for $(a,b,b)\in C_0$; replace two letters of a linear set by the third, or
delete the three letters of a linear set, in any order (an order that is not a cyclic rotation of the listed one costs three uses of (E1)
more).
\end{enumerate}
Each step multiplies the word by a conjugate of one relator and a bounded number of relators $x^2$, so counting relator applications bounds
the area. Each occurrence of $[p,\iota p\iota]$ is charged the area of the s-lemma, whose resolved sequence is
$\iota_1\,\iota_2\,\iota_1\,\iota\,\sigma$: one move of $C_0$ gives $(\iota_1\iota_2\iota_1)\,\iota\,\sigma$, and one linear relation deletes
it. The last derivation is $xzxz\,w\mapsto(xzx)\,z\,w\mapsto e$, by the triple $(x,z,xzx)$ and the linear set $\{xzx,z,w\}$. An exact count
of cells in the free group on $S_0\sqcup S_1$, charging every relator and every use of (E1), gives area at most $60$ for each of the $29$
words, attained by $(sps\,p^{-1})^3$: six uses of the s-lemma at $9$ cells each, and $6$ more (\cref{cert:what}). Two programs that share
no code establish this bound, one by a fresh search and one by writing each recorded derivation out as an explicit product of conjugates of
relators and checking it by free reduction. A third re-justifies every move as a rewriting along a cyclic rotation of a relator and
confirms the derivations with a more generous count, at most $335$ cells per word; this weaker bound would serve equally well, as $A_*$
enters only constants. No derivation uses the $25$
triples $(g,w,w)$ or the $19$ added commutations.
\end{proof}

Any derivation can be replayed by hand. For example, $qxq^{-1}x^{-1}$ becomes $\iota p\iota\,x\,\iota p^{-1}\iota\,x^{-1}$, with resolved
sequence $\iota\,\iota_1\,x\,\iota_1\,\iota\,x$ (the factor $pxp^{-1}$ resolves to $x$ by $C_1$). The triple $(\iota_1,x,x)$ turns it into
$\iota\,x\,\iota\,x$, the triple $(x,\iota,\iota)$ into $\iota\,\iota$, and (E1) cancels it.

\paragraph{What the presentation is for.} Three properties of $\Pres$ and $w$ matter. The first two are the hypotheses of the embedding
theorem of Slofstra and Vidick~\cite[Prop.~6.2]{SV18}; the third ties $\Pres$ to the certificate.
\begin{enumerate}[label=(R\arabic*),itemsep=1pt]
\item \emph{$w$ is a central involution} of $\langle\Pres\rangle$: it commutes with every generator by \cref{prop:presentation}(c), and
$w^2=e$ by (E1). Its image $\pi(w)=J$ is the central involution of $\GPp$.
\item \emph{$w$ survives}: for every $\eps>0$, $\Pres$ has a finite-dimensional $\eps$-representation with $w\mapsto-I$ (\cref{game:models}
below), so $w$ is nontrivial in approximate representations. This comes from the amenability of $\GPp$ through its hash models, and gives
the upper bound and the value one.
\item \emph{$\RP$ follows from $\Pres$}, with area at most $A_*=60$ (\cref{prop:presentation}(d)). This gives the lower bound.
\end{enumerate}
Property (R1) is what the flat embedding theorem of Slofstra and Vidick needs: there a visible central involution is amplified to
$-I$~\cite[Prop.~2.7]{SV18}. The state-weighted proof below never amplifies, and uses only (R2) and (R3).

\begin{lemma}[$w$ survives]\label{game:models}
For every $0<\eps\le1$ there is an $\eps$-representation $\phi$ of $\Pres$ with $\phi(w)=-I$ in a dimension $D$ with
\[
\log_2D\ \le\ 400\,\ell_E^2\,\eps^{-2}\log_2(2\ell_E/\eps),\qquad\ell_E=44 .
\]
\end{lemma}

\begin{proof}
Let $\vartheta:F(S_0\sqcup S_1)\to F(p,\iota,t,d,x,z)$ be the homomorphism of \cref{prop:presentation}(b), and let $\mathcal W$ be the
finite set of the free reductions of the words $\vartheta(r)$, $r\in R$. They represent $1$ in $\GPp$ and have length at most $\ell_E$. By
\cref{prop:hash-models} there are $D$ as stated and $\phi_0:F(p,q,\iota,t,d,x,z)\to\U(D)$ with $\fn{\phi_0(u)-I}\le\eps$ for $u\in\mathcal W$
and $\phi_0((xz)^2)=-I$. Since $\phi_0(\vartheta(r))$ depends only on the free reduction of $\vartheta(r)$, $\phi=\phi_0\circ\vartheta$ is an
$\eps$-representation of $\Pres$, and $\phi(w)=\phi_0(xzxz)=-I$.
\end{proof}

\subsection{Theorem A}

Let $(A,b)$ be the binary linear system that \cref{lem:gadget} builds from $\Pres$ and $w$. It is listed in \texttt{anc/game\_Ab.json};
\cref{box:game} describes it and \cref{rem:game-size} counts its parts. Let $G=G(A,b)$ be its game with the default question distribution
of \cref{def:lcs}, and $\Gamma=\Gamma(A,b)$ its solution group. Recall $\bfun(\eps)=2^{-4(24+\log\log(1/\eps))^2}$, which is
nondecreasing on $(0,\frac12)$ and equals $\eps^{o(1)}$ as $\eps\to0$.

\begin{figure}[t]
\centering
\fbox{\begin{minipage}{0.95\textwidth}\small
\textbf{The game at a glance.} We write $x_a$ for the variable of a generator $a$.
\begin{itemize}[leftmargin=1.2em,itemsep=1pt,topsep=2pt]
\item \emph{Size.} $m=1687$ equations over $\F_2$ in $n=1993$ variables, each equation in at most $s=3$ variables. The system is in
\texttt{anc/game\_Ab.json}; it is generated from $\Pres$ and checked by machine (\cref{rem:game-size,app:cert}).
\item \emph{Questions.} The referee picks an equation $i$ uniformly at random and then one of its variables $j$ uniformly at random.
Alice receives $i$ and Bob receives $j$.
\item \emph{Answers.} Alice answers values of the variables of equation $i$ that satisfy it, one of $4$ answers for three variables, and
Bob answers one bit. They win if they give $x_j$ the same value.
\item \emph{Equations.} All are homogeneous except $x_w+x_u=1$. There are four kinds (\cref{lem:gadget}): the $8$ linear relations of
$\Pres$, such as $x_{xzx}+x_z+x_w=0$, where $w$ stands for $xzxz$; the equation $x_w+x_u=1$, which makes $J$ a commutator; $274$ equations
in three variables, such as $x_a+x_{a'}+x_{a''}=0$, that put the $97$ conjugacy relations in a special form; and, for each of the
resulting $234$ relations $aba=c$, six equations in seven new variables,
$x_a+y_1+y_2=x_b+y_2+y_3=y_3+y_4+y_5=x_a+y_5+y_6=x_c+y_6+y_7=y_1+y_4+y_7=0$.
\item \emph{Value.} $\omq(G)=1$, but no strategy on finite-dimensional spaces wins with certainty, and hence none on a tensor product of
Hilbert spaces of any dimension~\cite[Thm.~3.2]{Slo19}. The system has no solution over $\F_2$.
\item \emph{Entanglement.} Every strategy that wins with probability $1-\eps$ has entanglement entropy at least $\eps^{-1+o(1)}$, and
a maximally entangled strategy with $\eps^{-1+o(1)}$ ebits suffices (\cref{thm:main}).
\end{itemize}
\end{minipage}}
\caption{The game $G$ of Theorem~A.}
\label{box:game}
\end{figure}

\begin{theorem}[Theorem~A]\label{thm:main}
The game $G$ has $m=1687$ equations in $n=1993$ variables, with at most $s=3$ variables in each equation, and:
\begin{enumerate}[label=(\alph*),itemsep=1pt]
\item $\omq(G)=1$, and no strategy on finite-dimensional spaces wins $G$ with probability one; hence, by~\cite[Thm.~3.2]{Slo19}, no
strategy on a tensor product of Hilbert spaces of any dimension does.
\item Let $c=4.9\cdot10^{-34}$ and $\eps_0=2^{-5500}$. For $0<\eps\le\eps_0$, every strategy that wins $G$ with probability at least $1-\eps$
satisfies
\[
\log\rank\rho_B\ \ge\ S(\rho_B)\ \ge\ c\,\eps^{-1}\bfun(\eps),
\]
where $\rho_B$ is Bob's reduced density matrix.
\item For $0<\eps\le\eps_0/2$, every mixed state $\sigma$ and measurements that win $G$ with probability at least $1-\eps$ satisfy
$E_F(\sigma)\ge\frac c4\,\eps^{-1}\bfun(\eps)$ (\cref{cor:main-EF}).
\item For $0<\eps\le\frac12$, $\Entme(G,\eps)\le\exp\bigl(C\eps^{-1}\log(1/\eps)\bigr)$ with $C=7.4\cdot10^{10}$.
\end{enumerate}
Consequently $\log\Ent(G,\eps)=\eps^{-1+o(1)}=\log\Entme(G,\eps)$ as $\eps\to0$.
\end{theorem}

All constants are explicit and far from optimal. The constants $c$ and $\eps_0$ come from the constant $A_0=42$ of \cref{lem:gadget}, the
largest area in $\Gamma$ of finitely many words that the construction makes trivial (\cref{rem:A0}), and from the sector theorem
(\cref{game:constants}).

\subsection{The proof at a glance}

\Cref{game:fig-map} shows the two chains of the proof.

\emph{The lower bound.} A strategy that wins with probability $1-\eps$ gives, on its Schmidt support, a state-weighted model $(\phi,\rho)$
of the solution group $\Gamma$ with $\phi(J)=-I$ (\cref{thm:transfer}): every relator has defect at most $4\theta$ in Bob's reduced state
$\rho$, and every generator nearly commutes with $\xi=\rho^{1/2}$, where $\theta=2\sqrt{ms\eps}$. So the precision is of order $\sqrt\eps$, and it
is never flattened. The substitution $\Psi$ of the embedding pulls the model back to a model of $\Pres$ at a constant factor worse
precision, and the commutator relation $[x_w,x_{j_0}]J=e$ built into the game forces $\phi(x_w)$ to stay at distance about $1$ from $I$ in
the state (\cref{thm:embedding}(b)): this is where $J=-I$ becomes visible on the side of $\Pres$. The words $\hat a$ of
\cref{prop:presentation}(d) pull the model back once more, to an assignment of unitaries to $p,q,t,d,x,z$ that satisfies $\RP$ in the state
and in which $J=(xz)^2$ stays visible. Then the sector theorem counts approximate Pauli pairs on the sector where $J$ acts as $-1$, and
returns entropy $e^{-2}$, up to the factor $\bfun$, at precision $e$ (\cref{thm:sector}); \cref{lem:sector-entropy} moves the entropy from
the sector state to $\rho$. Every step costs a constant factor in the precision, so $e^{-2}\asymp(\sqrt\eps)^{-2}=\eps^{-1}$ survives.

\emph{The upper bound.} The hash models give approximate representations of $\GPp$ with $J\mapsto-I$ exactly, of defect $\eps'$ in
dimension $\exp(O(\eps'^{-2}\log(1/\eps')))$ (\cref{prop:hash-models}). Composed with the defining words they are $\eps'$-representations of
$\Pres$ with $w\mapsto-I$ (\cref{game:models}). Slofstra's gadgets turn these into $3\eps'$-representations of $\Gamma$ with $J\mapsto-I$, in
eight times the dimension (\cref{thm:embedding}(a)). With the flat state, the return direction of the transfer turns these into strategies
on a maximally entangled state that lose with probability $O(\eps'^2)$ (\cref{thm:return}). Taking $\eps'\asymp\sqrt\eps$ gives local
dimension $\exp(O(\eps^{-1}\log(1/\eps)))$.

\begin{figure}[t]
\centering
\begin{tikzpicture}[
  st/.style={draw,rounded corners=2pt,align=flush center,font=\footnotesize,text width=4.9cm,inner sep=3pt},
  ar/.style={-{Stealth[length=2mm]},semithick},
  lb/.style={font=\scriptsize,align=flush left,anchor=west,text width=2.8cm}]
\node[font=\small\bfseries] at (0,0.8) {lower bound};
\node[st] (l1) at (0,0) {a strategy that wins with\\ probability $1-\eps$};
\node[st] (l2) at (0,-1.75) {a model $(\phi,\rho)$ of $\Gamma$, $\phi(J)=-I$,\\ precision $e_\Gamma=4\theta$, $\theta=2\sqrt{ms\eps}$};
\node[st] (l3) at (0,-3.5) {a model of $\Pres$ at precision $e_\Pres$,\\ $\hs{(\phi(x_w)-I)\xi}\ge1-e_\Pres$};
\node[st] (l4) at (0,-5.25) {a model of $\RP$ at precision $e$,\\ $\hs{(\phi'(J)-I)\xi}\ge\frac12$};
\node[st] (l5) at (0,-7.0) {$S(\rho_-)\gtrsim e^{-2}\bfun(e)$, so\\ $S(\rho_B)\ge c\,\eps^{-1}\bfun(\eps)$};
\draw[ar] (l1)--node[lb,right]{\cref{thm:transfer}}(l2);
\draw[ar] (l2)--node[lb,right]{$\Psi$; \cref{thm:embedding}(b)}(l3);
\draw[ar] (l3)--node[lb,right]{the words $\hat a$;\\ \cref{prop:presentation}(d)}(l4);
\draw[ar] (l4)--node[lb,right]{\cref{thm:sector};\\ \cref{lem:sector-entropy}}(l5);
\node[font=\small\bfseries] at (8.5,0.8) {upper bound};
\node[st] (u1) at (8.5,0) {hash model of $\GPp$, defect $\eps'$,\\ $\phi_0((xz)^2)=-I$};
\node[st] (u2) at (8.5,-1.75) {$\eps'$-representation of $\Pres$,\\ $w\mapsto-I$, dimension $D$};
\node[st] (u3) at (8.5,-3.5) {$3\eps'$-representation of $\Gamma$,\\ $J\mapsto-I$, dimension $8D$};
\node[st] (u4) at (8.5,-5.25) {strategy on $\Phi_{8D}$ losing with\\ probability at most $9C_3\eps'^2$};
\draw[ar] (u1)--node[lb,right]{$\vartheta$; \cref{game:models}}(u2);
\draw[ar] (u2)--node[lb,right]{\cref{thm:embedding}(a)}(u3);
\draw[ar] (u3)--node[lb,right]{\cref{thm:return},\\ flat state}(u4);
\end{tikzpicture}
\caption{The two directions of the proof of Theorem~A. Here $e_\Pres=C_\Psi e_\Gamma$ and $e=36780\,e_\Pres$; on the right,
$\eps'=\frac13\sqrt{\eps/C_3}$ with the constant $C_3<788.2$ of \cref{thm:return}.}
\label{game:fig-map}
\end{figure}

\subsection{Proof of the lower bound}

Throughout, $m=1687$ and $s=3$; the default distribution has $\pi_{\min}=1/(ms)$. Let $\Psi$, $x_{j_0}$ and $A_0=42$ be as in
\cref{lem:gadget} and \cref{rem:A0}, and $C_\Psi=A_0(8A_0+17)=14826$ as in \cref{thm:embedding}. Put
\[
C_*=294240\,\sqrt{ms}\;C_\Psi<3.11\cdot10^{11},\qquad c_*=\frac{c_1}{512\,C_*^2}>4.9\cdot10^{-34},\qquad c_1=\frac1{4.1\cdot10^7},
\]
where $c_1$ is the constant of \cref{thm:sector}. We prove (b) with $c_*$ in place of $c$, which is stronger, and for every
$\eps_0\in(0,\frac14]$ such that, for $0<\eps\le\eps_0$,
\begin{equation}\label{game:eps0}
C_*\sqrt\eps\le2^{-1072}\qquad\text{and}\qquad c_*\,\eps^{-1}\bfun(\eps)\ge1 ;
\end{equation}
such an $\eps_0$ exists because $\eps^{-1}\bfun(\eps)\to\infty$ as $\eps\to0$. The value $\eps_0=2^{-5500}$ of \cref{thm:main} qualifies.
Write $L=\log(1/\eps)\ge5500$. The first condition says $L\ge2144+2\log C_*$, and $2144+2\log C_*<2221$. The second says
$L-4(24+\log L)^2\ge\log(1/c_*)$, where $\log(1/c_*)<111$. Its left side is increasing for $L\ge5500$ and exceeds $190$ at $L=5500$.

\begin{proof}[Proof of \cref{thm:main}(b)]
Let $0<\eps\le\eps_0$, fix a strategy that wins with probability at least $1-\eps$, and put $\theta=2\sqrt{\eps/\pi_{\min}}=2\sqrt{ms\eps}$.

\emph{Step 1: the transfer.} By \cref{thm:transfer} there are a homomorphism $\phi:F(x_1,\dots,x_n,J)\to\U(r)$, where $r$ is the Schmidt
rank of the state, and a density matrix $\rho$ of full rank on $\C^r$, Bob's reduced density matrix restricted to its support, such that,
with $\xi=\rho^{1/2}$, $\phi(J)=-I$, every relator $R'$ of $\Gamma$ has $\hs{(\phi(R')-I)\xi}\le\max(s,4)\theta=4\theta$, and every $x_j$ has
$\hs{[\phi(x_j),\xi]}\le\sqrt2\theta$. Since $\phi(J)=-I$ commutes with $\xi$, $(\phi,\rho)$ is a state-weighted model of $\Gamma$
(\cref{def:model}) at precision $e_\Gamma=4\theta$.

\emph{Step 2: from $\Gamma$ to $\Pres$.} By \cref{thm:embedding}(b), $(\phi_\Pres,\rho)$ with $\phi_\Pres=\phi\circ\Psi$ is a state-weighted model
of $\Pres$ at precision $e_\Pres=C_\Psi e_\Gamma$, and $W=\phi_\Pres(w)=\phi(x_w)$ satisfies
\[
\hs{(W-I)\xi}\ \ge\ 1-e_\Pres .
\]

\emph{Step 3: from $\Pres$ to $\RP$.} Put $\phi'(a)=\phi_\Pres(\hat a)$ for $a\in\{p,q,t,d,x,z\}$, so that $\phi'(r)=\phi_\Pres(\hat r)$ for every
word $r$ in these letters. Let $r\in\RP$. It has at most $22$ letters (\cref{def:RP}), so $\abs{\hat r}\le66$, and $\hat r$ has area at most
$60$ with respect to $R$ by \cref{prop:presentation}(d). The relators of $\Pres$ have length at most $4$, so by \cref{lem:conjugator-length}
the word $\hat r$ is freely equal to a product of at most $60$ conjugates of relators of $\Pres$ and their inverses, with conjugators of length
at most $66+60\cdot4=306$. By \cref{lem:calculus}(d), with \ref{prelim:S1} and \ref{prelim:S2} for $\Pres$ at precision $e_\Pres$,
\[
\hs{(\phi'(r)-I)\xi}\ \le\ 60\,(1+2\cdot306)\,e_\Pres\ =\ 36780\,e_\Pres .
\]
The same holds for the word $xzxz\,w^{-1}$, which has $5$ letters and area at most $60$. Each $\hat a$ has at most $3$ letters, so
$\hs{[\phi'(a),\xi]}\le3e_\Pres$ by \cref{lem:calculus}(a). Put
\[
e=36780\,e_\Pres=147120\,C_\Psi\,\theta=C_*\sqrt\eps .
\]
Then $(\phi',\rho)$ satisfies conditions (i) and (ii) of \cref{thm:sector} at precision $e$, and $e\le2^{-1072}$ by~\eqref{game:eps0}.

\emph{Step 4: visibility.} Here $J$ denotes the word $xzxz$, so $\phi'(J)=\phi_\Pres(xzxz)$ and $\phi'(J)W^*=\phi_\Pres(xzxz\,w^{-1})$.
Since $\phi'(J)-W=(\phi'(J)W^*-I)W$ and $W\xi=\xi W+[W,\xi]$,
\[
\hs{(\phi'(J)-W)\xi}\ \le\ \hs{(\phi'(J)W^*-I)\xi}+2\hs{[W,\xi]}\ \le\ e+2e_\Pres ,
\]
using Step~3 and \ref{prelim:S2} for the generator $w$. With Step~2 and $3e_\Pres\le e\le\frac14$,
\[
\hs{(\phi'(J)-I)\xi}\ \ge\ \hs{(W-I)\xi}-\hs{(\phi'(J)-W)\xi}\ \ge\ 1-e-3e_\Pres\ \ge\ \tfrac12 .
\]

\emph{Step 5: the sector.} By Steps~3 and~4, \cref{thm:sector} applies with $\Delta=\frac12$, since $e\le2^{-1072}=2^{-1071}\Delta$. Let
$J'=\mathrm r(\phi'(J))$ be the rounding of \cref{lem:rounding}, $P_-=\frac12(I-J')$, $p_-=\Tr(P_-\rho)$, and $\rho_-$ the sector state of
\cref{lem:sector-state} for $J'$. By \cref{thm:sector}, $p_-\ge\Delta^2/16=\frac1{64}$ and
\[
S(\rho_-)\ \ge\ \tfrac{c_1}4\,e^{-2}\,\bfun(e)\ \ge\ \frac{c_1}{4C_*^2}\,\eps^{-1}\bfun(\eps),
\]
because $e^{-2}=C_*^{-2}\eps^{-1}$, and $\eps\le\sqrt\eps\le e<\frac12$ (as $C_*\ge1$) with $\bfun$ nondecreasing on $(0,\frac12)$.

\emph{Step 6: back to $\rho$.} By \cref{lem:sector-entropy}, since $h\le1$,
\[
S(\rho)\ \ge\ p_-\,S(\rho_-)-1\ \ge\ \frac{c_1}{256\,C_*^2}\,\eps^{-1}\bfun(\eps)-1\ =\ 2c_*\,\eps^{-1}\bfun(\eps)-1\ \ge\ c_*\,\eps^{-1}\bfun(\eps),
\]
the last step by~\eqref{game:eps0}. Finally $\rho$ and $\rho_B$ have the same nonzero eigenvalues, so
$\log\rank\rho_B=\log\rank\rho\ge S(\rho)=S(\rho_B)$.
\end{proof}

The proof uses the strategy only through \cref{thm:transfer}, and the state only through \ref{prelim:S1} and \ref{prelim:S2}; no spectral
projection of $\rho$ is taken. Since $\omq(G)=1$ by (a), a strategy on $\C^d\otimes\C^d$ that wins with probability at least
$\omq(G)-\eps$ has Schmidt rank at most $d$, and (b) applies to it; so $\log\Ent(G,\eps)\ge c\,\eps^{-1}\bfun(\eps)$ for $0<\eps\le\eps_0$.

\begin{corollary}[Entanglement of formation]\label{cor:main-EF}
With $c$ and $\eps_0$ as in \cref{thm:main}(b), let $0<\eps\le\eps_0/2$. Every mixed state $\sigma$ on $H_A\otimes H_B$ and measurements
that win $G$ with probability at least $1-\eps$ satisfy
\[
E_F(\sigma)\ \ge\ \tfrac c4\,\eps^{-1}\bfun(\eps).
\]
\end{corollary}

\begin{proof}
Let $\sigma=\sum_jp_j\lvert\psi_j\rangle\langle\psi_j\rvert$ with $p_j>0$ and unit vectors $\psi_j$, and let $\ell_j$ be the losing
probability of the pure strategy that uses $\psi_j$ with the given measurements. The losing probability is linear in the state, so
$\sum_jp_j\ell_j\le\eps$, and the components with $\ell_j>2\eps$ carry total weight less than $\frac12$. Each component with $\ell_j\le2\eps$
is a strategy that wins with probability at least $1-2\eps$, with $2\eps\le\eps_0$, so by \cref{thm:main}(b) and the monotonicity of
$\bfun$,
\[
S\bigl(\Tr_A\lvert\psi_j\rangle\langle\psi_j\rvert\bigr)\ \ge\ c\,(2\eps)^{-1}\bfun(2\eps)\ \ge\ \tfrac c2\,\eps^{-1}\bfun(\eps).
\]
These components carry weight at least $\frac12$, so $\sum_jp_jS(\Tr_A\lvert\psi_j\rangle\langle\psi_j\rvert)\ge\frac c4\eps^{-1}\bfun(\eps)$.
The infimum over decompositions is $E_F(\sigma)$ (\cref{def:entanglement}).
\end{proof}

\subsection{Proof of the upper bound and of value one}

\begin{proof}[Proof of \cref{thm:main}(d) and (a)]
Let $C_3<788.2$ be the constant $C_s$ of \cref{thm:return} for $s=3$. Let $0<\eps\le\frac12$ and put $\eps'=\frac13\sqrt{\eps/C_3}\le1$.

\emph{Step 1: a representation of $\Pres$.} By \cref{game:models} there is an $\eps'$-representation of $\Pres$ with $w\mapsto-I$ in a
dimension $D$ with $\log_2D\le400\cdot44^2\,\eps'^{-2}\log_2(88/\eps')$.

\emph{Step 2: a representation of $\Gamma$.} By \cref{thm:embedding}(a) there is a $3\eps'$-representation $\phi_\Gamma$ of $\Gamma$ on
$\C^{8D}$ with $\phi_\Gamma(J)=-I$.

\emph{Step 3: a maximally entangled strategy.} Let $\rho=I/(8D)$ be the flat state. Then $\hs{Y\xi}=\fn{Y}$ for every operator $Y$, all state
commutators vanish, and $\phi_\Gamma(J)+I=0$, so $(\phi_\Gamma,\rho)$ is a model of $(\Gamma,J)$ at precision $3\eps'$
(\cref{def:model}). By \cref{thm:return}, the strategy built there on the canonical purification of $\rho$, which is the maximally
entangled state $\Phi_{8D}$, has local dimension $8D$ and loses every question pair, hence a random one, with probability at most
$C_3(3\eps')^2=\eps$.

\emph{Step 4: the count.} Since $\eps'^{-2}=9C_3\eps^{-1}$ and $88/\eps'=264\sqrt{C_3/\eps}$,
\[
\log_2(8D)\ \le\ 3+400\cdot44^2\cdot9\,C_3\,\eps^{-1}\Bigl(\log_2\bigl(264\sqrt{C_3}\bigr)+\tfrac12\log_2(1/\eps)\Bigr).
\]
Here $400\cdot44^2\cdot9\,C_3<5.50\cdot10^9$ and $\log_2(264\sqrt{C_3})<12.86$, and for $\eps\le\frac12$ we have $\log_2(1/\eps)\ge1$ and
$3\le\frac32\eps^{-1}\log_2(1/\eps)$. So $\log_2(8D)\le7.4\cdot10^{10}\,\eps^{-1}\log_2(1/\eps)$. Since $2^y\le e^y$ for $y\ge0$, the strategy
of Step~3 shows $\Entme(G,\eps)\le\exp(C\eps^{-1}\log(1/\eps))$ with $C=7.4\cdot10^{10}$, once we know $\omq(G)=1$.

\emph{Step 5: value one.} Step~3 gives, for every $0<\eps\le\frac12$, a strategy that wins with probability at least $1-\eps$. So
$\omq(G)=1$, which completes (d). Suppose a strategy on finite-dimensional spaces won with probability one. It would win with probability at
least $1-\eps$ for every $\eps$, so by (b) its reduced density matrix would satisfy $\log\rank\rho_B\ge c\,\eps^{-1}\bfun(\eps)$ for every
$0<\eps\le\eps_0$. The right side is unbounded as $\eps\to0$, while $\rank\rho_B\le\dim H_B$. So there is no such strategy. By
\cite[Thm.~3.2]{Slo19}, a linear system game has a perfect strategy on a tensor product of Hilbert spaces of arbitrary dimension if and only
if it has one on finite-dimensional spaces; so there is none of that kind either.
\end{proof}

\begin{proof}[Proof of the last sentence of \cref{thm:main}]
For $0<\eps\le\eps_0$, the paragraph after the proof of (b), $\Ent\le\Entme$ and (d) give
\[
c\,\eps^{-1}\bfun(\eps)\ \le\ \log\Ent(G,\eps)\ \le\ \log\Entme(G,\eps)\ \le\ C\log_2(e)\,\eps^{-1}\log(1/\eps).
\]
Since $\bfun(\eps)=\eps^{o(1)}$ and $\log(1/\eps)=\eps^{-o(1)}$, both logarithms are $\eps^{-1+o(1)}$.
\end{proof}

\subsection{Remarks}

\begin{remark}[The size of the game]\label{rem:game-size}
The four steps of \cref{lem:gadget}, applied to $\Pres$, give the following counts, which the generating program and the independent
checker of \cref{app:cert} both reproduce. Step~1 adds $12$ involutions ($f$, $v$, and one $Z_{bc}$ for each of the $10$ triples of $C_1$)
and $20$ conjugacy relations to the $76$ of $\Pres$. Step~2 adds the involutions $t^*$, $u$, the linear set $\{w,u\}$ and one conjugacy
relation; this leaves $N=40$ generators besides $J$ and $97$ conjugacy relations. Step~3 adds $3N+1+2\cdot97=315$ involutions and
$2N+2\cdot97=274$ linear sets, and turns the $97$ conjugacy relations into $N+2\cdot97=234$ nice ones. Step~4 replaces each nice relation
by $6$ equations in $7$ new variables. So $8+1+274=283$ equations come from linear sets and $6\cdot234=1404$ from gadgets, $1687$ in all,
in $26+12+2+315+1638=1993$ variables. All equations have three variables except $x_w+x_u=1$, which has two and is the only inhomogeneous
one. Gaussian elimination shows that $Ax=b$ has no solution over $\F_2$ ($A$ has rank $1512$), as it must, since $J$ is a commutator in
$\Gamma$: no deterministic strategy wins with certainty. A program independent of the generator also evaluates every relator of $\Gamma$ in
the representations $\gamma,\gamma_1,\gamma_2$ of the proof of \cref{lem:gadget}, built from the exact model of $\GPp$ with its central
involution $J$ in place of $-1$, and finds a homomorphism $\Gamma\to\GPp\wr\Sym(8)$ in which $J\ne1$ (\cref{cert:what}). None of this is used
in the proof of \cref{thm:main}, which needs only that $(A,b)$ is the output of \cref{lem:gadget}.
\end{remark}

\begin{remark}[The constants]\label{game:constants}
Only the constants $c$ and $\eps_0$ depend on the area constant $A_0=42$ of \cref{rem:A0}, and they are astronomically small: already
the precision condition $e\le2^{-1071}\Delta$ of the sector theorem forces $\eps_0<2^{-2144}$, and the factor $\bfun(\eps)$ pushes the
admissible $\eps_0$ down to about $2^{-5420}$. The constant $C$ of (d) is explicit, and
improving it is a matter of the hash models (the factor $400\,\ell_E^2$) and of the constant of \cref{thm:return}. The bounds (b) and (d) have
the same power of $\eps$; only the factors $\bfun(\eps)$ and $\log(1/\eps)$ separate them.
\end{remark}

\begin{remark}[Under the polynomial Dehn bound]\label{game:poly}
If the polynomial bound~\ref{H-poly} on the Dehn function of $H_3$ is used, \cref{rem:sector-poly} replaces the factor
$2^{-4(24+\log\log(1/e))^2}$ of \cref{thm:sector} by a constant multiple of $(\log(1/e))^{-4\kappa}$, for $e\le e_0\Delta$. The proof above
then gives \cref{thm:main}(b),(c) and \cref{cor:main-EF} with $\bfun(\eps)$ replaced by $(\log(1/\eps))^{-4\kappa}$, with new constants
$c,\eps_0$, since $e\ge\eps$ gives $(\log(1/e))^{-4\kappa}\ge(\log(1/\eps))^{-4\kappa}$.
\end{remark}
